{\large\textbf{Space Debris}}

\textbf{Introduction}

In more than half a century of space operations quite a large number of
man-made objects have been amassed near Earth. The objects that do not serve
any particular purpose are called \textbf{space debris}. The most attention is
usually paid to the larger debris objects, i.e. defunct satellites and spent
rocket upper stages, which stay in orbit after delivering their payload.
Collisions of such objects with each other may result in thousands of fragments
endangering all current space missions.

There is a well-known hypothetical scenario, according to which certain
collisions may cause a cascade where each subsequent collision generates more
space debris that increase the likelihood of new collisions. Such a chain
reaction, resulting in the loss of all near-Earth satellites and making
impossible further space programs, is called the \textbf{Kessler syndrome}.

To prevent such undesirable outcome special missions are planned to remove
large debris object from their present orbits either by tugging them to the
Earth's atmosphere or to graveyard orbits. To this end a specially designed
spacecraft – a space tug – must capture a debris object. However, before
capturing an uncontrolled object it is important to understand its rotational
dynamics.

We suggest you to take part in planning of such a mission and find out how the
rotational dynamics of a debris object changes in time under the influence of
different factors.

\textbf{Rocket Stage Schematic}

The debris object to be considered is a \textquotedbl Kerbodyne 42\textquotedbl{} rocket upper stage,
whose schematic is shown in Fig. 1. The circle line in Fig. 1 marks the outline
of a spherical fuel tank.

\begin{center}\includegraphics[width=0.48\linewidth]{figures/APhO_2017_Q3_fig1.png}\end{center}

\begin{center}Fig. 1: \textquotedbl Kerbodyne 42\textquotedbl{} upper stage\end{center}

We introduce a body-fixed reference frame $Cxy$ with the origin in the center
of mass $C$, $x$ being the symmetry axis of the stage, and $y$ perpendicular to
$x$. The inertia moments with respect to $x$ and $y$ axes are $J_x$ and $J_y$
($J_x < J_y$).

\textbf{Part A. Rotation (3.8 points).}

Consider an arbitrary initial rotation of the stage with angular momentum $L$
(Fig. 2), where $\theta$ is the angle between the symmetry axis and the
direction of angular momentum. Fuel tank at this point is assumed to be empty.
No forces or torques act upon the stage.

\begin{center}\includegraphics[width=0.46\linewidth]{figures/APhO_2017_Q3_fig2.png}\end{center}

\begin{center}Fig. 2: Rocket stage rotation\end{center}

\textbf{A.1}\quad Find the projections of angular velocity $\vec{\omega}$ on
$x$ and $y$, given that
$\vec{L} = J_x\omega_x\vec{e}_x + J_y\omega_y\vec{e}_y$ for material symmetry
axes $x$ and $y$ with unit vectors $\vec{e}_x$ and $\vec{e}_y$. Provide the
answer in terms of $L = |\vec{L}|$, angle $\theta$, and inertia moments $J_x$,
$J_y$. \hfill 0.2pt

\textbf{A.2}\quad Find the rotational energy $E_x$ associated with rotation
$\omega_x$ and $E_y$ associated with rotation $\omega_y$. Find total rotational
kinetic energy $E = E_x + E_y$ of the stage as a function of the angular
momentum $L$ and $\cos\theta$. \hfill 0.4pt

In the following questions of Section A consider the stage’s free rotation
with the initial angular momentum $L$ and $\theta(0) = \theta_0$.

\textbf{A.3}\quad Let us denote by $x_0$ the initial orientation of the stage's
symmetry axis $Cx$ with respect to inertial reference frame. Using conservation
laws find the maximum angle $\psi$, which the stage's symmetry axis $Cx$ makes
with $x_0$ during the stage's free rotation.\\
\textit{Note:} Since there are no external torques acting upon the stage, the
angular momentum vector remains constant. \hfill 1.2pt

\begin{center}\includegraphics[width=0.47\linewidth]{figures/APhO_2017_Q3_fig3.png}\end{center}

\begin{center}Fig. 3: Precession\end{center}

Let us now introduce the reference frame $Cx_1y_1z_1$ with $y_1$ along the
constant angular momentum vector $\vec{L}$ (Fig. 3). This reference frame
rotates about $y_1$ in such a way, that the stage's symmetry axis always
belongs to the $Cx_1y_1$ plane.

\textbf{A.4}\quad Given $L$, $\theta(0) = \theta_0$, and inertia moments
$J_x, J_y$, find the angular velocity $\Omega(t)$ of the reference frame
$Cx_1y_1$ about $y_1$ and direction and absolute value of angular velocity of
the stage $\vec{\omega}_s(t)$ relative to the reference frame $Cx_1y_1$ as
functions of time. Provide the answer for $\vec{\omega}_s(t)$ direction in
terms of angle $\gamma_s(t)$ it makes with the symmetry axis $Cx$.\\
\textit{Note}: angular velocity vectors are additive
$\vec{\omega} = \vec{\omega}_x + \vec{\omega}_y = \vec{\Omega} + \vec{\omega}_s$.
\hfill 2.0pt

\textbf{Part B. Transient Process (1.6 points).}

Most of the propellant is used during the ascent, however, after the payload
has been separated from the stage, there still remains some fuel in its tank.
Mass $m$ of residual fuel is negligible in comparison to the stage's mass $M$.
Sloshing of the liquid fuel and viscous friction forces in the fuel tank result
in energy losses, and after a transient process of irregular dynamics the
energy reaches its minimum.

\textbf{B.1}\quad Find the value $\theta_2$ of angle $\theta$ after the
transient process for arbitrary initial values of $L$ and
$\theta(0) = \theta_1 \in (0, \pi/2)$ . \hfill 0.6pt

\begin{center}\includegraphics[width=0.45\linewidth]{figures/APhO_2017_Q3_fig4.png}\end{center}

\begin{center}Angle between the stage's angular velocity and the symmetry axis\end{center}

\textbf{B.2}\quad Calculate the value $\omega_2$ of angular velocity $\omega$
after the transient process, given that initial angular velocity
$\omega(0) = \omega_1 = 1\,rad/s$ makes an angle of
$\gamma(0) = \gamma_1 = 30^\circ$ with the stage's symmetry axis. The moments
of inertia are $J_x = 4200\ kg \cdot m^2$ and $J_y = 15\ 000\ kg \cdot m^2$.
\hfill 0.6pt

\textbf{Part C. Magnetic Field (4.6 points).}

Another important factor in rotational dynamics of a debris rocket stage, which
is orbiting the Earth, is its interaction with the Earth's magnetic field. Let
us first consider an auxiliary problem.

\textbf{Torque due to Eddy Currents.}

Let us place a thin-walled nonmagnetic spherical shell with wall thickness $D$
and radius $R$ in a uniform magnetic field $\vec{B}$, which slowly changes so
that its derivative $\dot{\vec{B}}$ is a constant vector making angle $\alpha$
with the direction of $\vec{B}$ (Fig. 4). Electrical resistivity of the
shell's material is $\rho$.

\begin{center}\includegraphics[width=0.47\linewidth]{figures/APhO_2017_Q3_fig5.png}\end{center}

\begin{center}Fig. 4: Spherical shell in magnetic field\end{center}

\textbf{C.1}\quad Find the induced magnetic moment $\vec{\mu}$ of the shell,
neglecting its self-inductance. Provide the answer for $\vec{\mu}$ in the form
of projections on $xyz$ (see Fig. 4). \hfill 1.0pt

\textbf{C.2}\quad Find the torque $\vec{M}$ acting on the spherical shell.
Provide the answer for $\vec{M}$ in the form of projections on $xyz$ (see
Fig. 4). \hfill 0.3pt

\textbf{Attitude Motion Evolution in the Earth’s Magnetic Field}

Let us find out how the rotation changes for a rocket stage, which moves in a
circular polar orbit with orbital period $T = 100\ min$ (Fig. 5). It transpires
that the characteristic times of dynamics due to interaction with the
geomagnetic field are much greater than the duration of the transient process.
We will now study what happens to the rocket stage after the transient process
has completed. To start our analysis consider the stage rotating with angular
velocity $\omega_2$ about the axis perpendicular to the orbital plane.

\begin{center}\includegraphics[width=0.50\linewidth]{figures/APhO_2017_Q3_fig6.png}\end{center}

\begin{center}Fig. 5: The orbit\end{center}

\textbf{C.3}\quad The Earth's magnetic field $\vec{B}_E$ can be modeled as the
magnetic field of a point dipole in the Earth's center. Its dipole moment
$\vec{\mu}_E$ is directed opposite to $Y$ axis. The absolute value of the
Earth's magnetic field $B$ at the point where the orbit crosses the equatorial
plane $XZ$ is $B_0 = 20\ \mu T$. Find $\vec{B}_E(u)$ at a current position of
the stage in the orbit defined by the angle $u$ as shown in Fig. 5. The
positive direction of $u$ is along with the orbital motion. Provide the answer
in the form of the projections of $\vec{B}_E(u)$ on $XYZ$.\\
\textit{Note}: It may facilitate subsequent calculations if projections of
$\vec{B}_E(u)$ are given as functions of $2u$ instead of $u$.\\
\textit{Note}: Magnetic field of a dipole at point $\vec{r}$ is given by
\begin{equation}
\vec{B} = \frac{\mu_0}{4\pi}\left(\frac{3(\vec{\mu}\cdot\vec{r})\vec{r}}{r^5} - \frac{\vec{\mu}}{r^3}\right).
\tag{1}
\end{equation}
\hfill 0.4pt

The \textquotedbl Kerbodyne 42\textquotedbl{} rocket upper stage is mostly made of wood, and the only
conductive material is used for its cryogenic fuel tank. We, therefore,
consider the stage’s interaction with the geomagnetic field as that of the
spherical shell with wall thickness $D = 2\,\mathrm{mm}$, radius
$R = 4\,\mathrm{m}$ and resistivity
$\rho = 2.7 \cdot 10^{-8}\,\Omega \cdot \mathrm{m}$.

\textbf{C.4}\quad Find the torque $\vec{M}(u)$ acting on the stage, as it
rotates about the axis perpendicular to the orbital plane with angular velocity
$\omega$ collinear to $Z$. Provide the answer in the form of the projections of
$\vec{M}(u)$ on $XYZ$. \hfill 1.3pt

\textbf{C.5}\quad Find the absolute value of angular velocity $\omega(t)$ as a
function of time, given that the change in the stage’s angular velocity over
one orbital period is negligibly small. \hfill 1.0pt

\textbf{C.6}\quad Find the ratio of the orbital period $T$ and the rocket
stage’s rotation period $T_s$ in the steady-state regime, which sets in after a
long time. \hfill 1.0pt
